% Zone „Das kennst du schon“ – aus bank/terme/zone.jsonl (zusammenbau v0.9)
\einheitenkopf[zone][Das kennst du schon]{Das kennst du schon}
\setcounter{aufgabe}{0}


% zone: terme-zone-f1-v1, terme-zone-f1-v2, terme-zone-f1-v3, terme-zone-f1-v4
\begin{kbaufgabe}{Ich kann negative Zahlen addieren und subtrahieren.}
\begin{kbblock}
\kbanweisung{Rechne.}
\lbpaar{\kbteil{a)}{$4 - 9 =$ \lbfeld}{}
}{\kbteil{b)}{$-3 + 8 =$ \lbfeld}{}}
\end{kbblock}
\begin{kbblock}
\lbpaar{\kbteil{c)}{$-2{,}5 + 1 =$ \lbfeld}{}
}{\kbteil{d)}{$-6 - 4 =$ \lbfeld}{}}
\item[]\lbhaengst{6}
\end{kbblock}
\end{kbaufgabe}


% zone: terme-zone-f2-v1, terme-zone-f2-v2, terme-zone-f2-v3, terme-zone-f2-v4
\begin{kbaufgabe}{Ich kann Zahlen mit Vorzeichen malnehmen.}
\begin{kbblock}
\kbanweisung{Rechne.}
\lbpaar{\kbteil{a)}{$(-5) \cdot 2 =$ \lbfeld}{}
}{\kbteil{b)}{$6 \cdot (-2) =$ \lbfeld}{}}
\end{kbblock}
\begin{kbblock}
\lbpaar{\kbteil{c)}{$(-0{,}5) \cdot 8 =$ \lbfeld}{}
}{\kbteil{d)}{$(-4) \cdot (-5) =$ \lbfeld}{}}
\item[]\lbhaengst{6}
\end{kbblock}
\end{kbaufgabe}


% zone: terme-zone-f3-v1, terme-zone-f3-v2, terme-zone-f3-v3, terme-zone-f3-v4
\begin{kbaufgabe}{Ich kann mit Dezimalzahlen und einfachen Brüchen rechnen.}
\begin{kbblock}
\kbanweisung{Rechne.}
\lbpaar{\kbteil{a)}{$1{,}5 + 2{,}5 =$ \lbfeld}{}
}{\kbteil{b)}{$\frac{1}{2} + \frac{1}{4} =$ \lbfeld}{}}
\end{kbblock}
\begin{kbblock}
\lbpaar{\kbteil{c)}{$0{,}5 - 1{,}5 =$ \lbfeld}{}
}{\kbteil{d)}{$0{,}9 + 0{,}3 =$ \lbfeld}{}}
\item[]\lbhaengst{6}
\end{kbblock}
\end{kbaufgabe}


% zone: terme-zone-f4-v1, terme-zone-f4-v2, terme-zone-f4-v3, terme-zone-f4-v4
\begin{kbaufgabe}{Ich kann Punkt vor Strich rechnen.}
\begin{kbblock}
\kbanweisung{Rechne.}
\lbpaar{\kbteil{a)}{$5 + 2 \cdot 4 =$ \lbfeld}{}
}{\kbteil{b)}{$10 + 6 \cdot 3 =$ \lbfeld}{}}
\end{kbblock}
\begin{kbblock}
\lbpaar{\kbteil{c)}{$6 - 4 \cdot 2{,}5 =$ \lbfeld}{}
}{\kbteil{d)}{$9 - 2 \cdot 3 =$ \lbfeld}{}}
\item[]\lbhaengst{15}
\end{kbblock}
\end{kbaufgabe}


% zonepaar: terme-zone-f1-v5, terme-zone-f1-v6
\begin{kbaufgabe}{Ich finde den Fehler beim Rechnen mit negativen Zahlen.}
\begin{kbblock}
\kbteil{a)}{Tim soll $-7 - 5$ berechnen. Er rechnet so: \rechnung{-7 - 5 &= -2} Finde den Fehler und rechne richtig.}{}
\item[]\kbraum{2}
\end{kbblock}
\begin{kbblock}
\kbanweisung{Rechne.}
\kbteil{b)}{$-8 - 6 =$ \lbfeld}{}
\item[]\lbhaengst{6}
\end{kbblock}
\end{kbaufgabe}
